\section{Trees}

A surprisingly large amount of the study of programming languages depends on trees, and more generally on inductively defined structures.
This is why functional languages like OCaml are well-suited for the implementation of programming languages.
So we begin our discussion with the humble notion of a tree.

If you've taken a course in discrete mathematics (and I pray that you have by the time you're reading this), you're likely familiar with the graph-theoretic definition of trees.

\begin{definition}
  A \textbf{tree} is an connected acyclic undirected graph.
  A \textbf{directed tree} is a directed graph whose underlying graph is a tree.
  A directed tree is \textbf{rooted} if there is a unique vertex with in-degree 0.
\end{definition}
%
In the setting of programming languages we'll be exclusively interested in \textit{nonempty rooted directed} trees so, oving forward, by \enquote{tree} we will always mean \enquote{nonempty rooted directed tree.}
These are also sometimes called \textbf{rose trees}.
To make this more explicit we'll often work with the \textit{inductive} definition of a rose tree.

\begin{definition}
  \label{def:rose-tree}
  A \textbf{tree} with values from $V$ is defined inductively as follows:
  \begin{itemize}
  \item if $v$ is a value from $V$ and $T_1, \dots, T_k$ are trees then so is $\tnode(v, T_1, \dots, T_k)$
  \end{itemize}
  Note, in particular, that $\tnode(v)$ a tree for any value $v$ from $V$.
  The \textbf{root} of a tree is defined as follows:
  \begin{align*}
    \troot {\tnode(v, T_1, \dots, T_k)} = v
  \end{align*}
  The \textbf{children} a tree are defined as follows:
  \begin{align*}
    \tchildren{\tnode(v, T_1, \dots, T_k)} = [T_1, \dots, T_k]
  \end{align*}
\end{definition}

\begin{exercise}
  Rewrite \Cref{def:rose-tree} in OCaml.
\end{exercise}

\begin{example}
  Example
\end{example}

\begin{example}
  \label{ex:decision-tree}
  We use decision trees to represent boolean functions.
  Note that rose trees are suitable for representing decision trees because there is no need for an empty tree structure.
  The decision tree for the boolean function $\mathsf{OR}(x_1, x_2, x_3) = x_1 \lor x_2 \lor x_3$ can be written:
  \begin{align*}
    \tnode(x_1 = 1, \tnode(1), \tnode(x_2 = 1, \tnode(1), \tnode(x_3)))
  \end{align*}
  In this scenario, values are taken from a set of possible queries to the inputs of the function.
  Informally, this decision tree expresses the logic that the value of $\mathsf{OR}(x_1, x_2, x_3)$ can be determined by checking each input individually and determining if it's 1.
\end{example}

%% \section{Visualizing Trees}

I expect that you're already familiar with basic ways of visualizing trees.
Here, for example, is a visualization of the tree from \Cref{ex:decision-tree}:

\begin{center}
  \begin{forest}
    for tree={
      minimum width=1in,
    }
    [{$x_1 = 1$}
      [$1$]
      [{$x_2 = 1$}
        [$1$]
        [$x_3$]
      ]
    ]
  \end{forest}
\end{center}

We will occassionally also use the following form of tree visualization.
Throughout, we'll call this \textbf{compact form}.
We won't use it too often, but it's useful when a tree it to wide to draw otherwise.

\begin{center}
  \begin{forest}
    for tree={
      grow'=0,
      child anchor=west,
      parent anchor=south,
      anchor=west,
      calign=first,
      edge path={
        \noexpand\path [draw, \forestoption{edge}]
        (!u.south west) +(5pt,0) |- (.child anchor)\forestoption{edge label};
      },
      before typesetting nodes={
        if n=1
        {insert before={[,phantom]}}
        {}
      },
      %% fit=band,
      before computing xy={l=20pt},
    }
    [{$\Gamma \vdash M : A$}
      [text1.1
        [text1.1.1]
        [text1.1.2]
        [text1.1.3]
      ]
      [text1.2
        [text1.2.1]
        [text1.2.2]
      ]
    ]
  \end{forest}
\end{center}
%
Once we have a fixed notion of trees, we can define any nonempty inductive structure.

\begin{example}
  The natural numbers can be thought of as trees with values from $\{\bullet\}$, where $\mathsf{zero}$ is $\tnode(\bullet)$ and $\mathsf{succ}(n) = \tnode(\bullet, n)$.
\end{example}

%% \section{Structural Recursion and Induction}

Ultimately, we'll be interested in \textit{proving} things about trees.
This will nearly always come down to structural induction on trees.

Part of what we emphasize here is that, everything is trees.
Once we have a way to talk about induction over trees, we can talk about proving things in general.

\section{Formal Systems}

The study of programming languages is ultimately the study of a class of \textbf{formal systems}.
The programming language itself is just an \textit{implementation} of a formal system.
As such, we need to concern ourselves at least briefly with these notions.

Let's paint a picture.
When we program we type a bunch of symbols into a computer which get saved to a file.
Depending on what language we're using (including what Human language we're assuming) there is a fixed number of characters we're allowed to type into a file.
We call this set of symbols the \textbf{alphabet}.
Of the possible orientations of symbols we can type into a file, only a handful of them are going to be well-formed programs.
Of the well-formed programs, only a handful those make sense as programs.
It doesn't, for example, make sense to add a string to an integer; in other words, only a handful of well-formed programs are well-typed.
And of the well-typed programs, only a handful of \textit{those} make sense as computations.
the program
\begin{ocaml}
  let rec f x = f x
\end{ocaml}
Make sense as an OCaml program in the sense that we can \textit{try} to compute its value, but we'll never succeed (it's an infinite loop).

\todo{A picture, nested squares}

We're interested in formal systems of this form.  In broad strokes, a formal system is made up of a formal language, determined by a \textbf{formal grammar} and one or more \textbf{deductive systems}.
At each level, we have a deductive system which describes a subsection of system.
If each step is \textit{decidable} or \textit{semi-decidable}, meaning we can determine which satisfy what, then we can implement it as a programming language.

\section{Inference Rules}

Deductive systems are formalizations of deductive reasoning over what are called \textbf{judgments}.
They are made up of \textbf{inference rules}, which can be though of as functions which define when we can derive a new judgment from a collection of old ones.

\begin{definition}
  A \textbf{inference rule} is...
\end{definition}

\begin{example}

\end{example}

When we define an inference rule, we're often defining a \textit{collection} of inference rules simultaneously.
This is done by the use of \textbf{meta-variables}.
A meta-variable stands for a syntactic construct whose form is determined by context.

One we have a notion of inference rule, we need the notion of a \textbf{derivation}.

\begin{definition}
A \textbf{derivation} is a tree with the following properties.
\end{definition}

\section{Derivations}
